\chapter{Lyubich renormalization and a priori bounds}

\section{Polynomial-like maps}

\begin{definition}[Polynomial-like map]
\label{def:polynomial-like}
A polynomial-like map of degree $d$ is a proper holomorphic map
\[
 g:U\longrightarrow V
\]
where $U,V\subset\C$ are simply connected domains with
$\overline U\subset V$ and $d\ge2$.  Its filled Julia set is
\[
 K(g):=\{z\in U:g^{\circ n}(z)\in U\text{ for every }n\in\N\}.
\]
\end{definition}

\begin{definition}[Quadratic renormalization]
\label{def:quadratic-renormalization}
A quadratic polynomial $f_c$ is renormalizable at depth $m$ if there are
domains $U\Subset V$ around the critical orbit and an integer $p>1$ such
that
\[
 f_c^{\circ p}:U\to V
\]
is polynomial-like of degree two.  The renormalization is the normalized
polynomial-like map obtained from this restriction and its affine marking.
\end{definition}

\begin{definition}[Renormalization tower]
\label{def:renormalization-tower}
A renormalization tower is a sequence
\[
 g_0,\ g_1,\ g_2,\ldots
\]
of marked polynomial-like maps together with renormalization relations
$g_n\rightsquigarrow g_{n+1}$.  The cumulative depth is
\[
 d_0=0,\qquad d_{n+1}=d_n+p_n,
\]
where $p_n$ is the period of the $n$-th renormalization step.
\end{definition}

The Lean structure \lean{RenormalizationTower} stores the sequence and the
nonempty relation at each step.  The theorem
\lean{RenormalizationTower.cumulativePeriod_monotone} proves monotonicity of
the cumulative depths without adding analytic assumptions.

\section{Principal nests and modulus}

\begin{definition}[Principal nest]
\label{def:principal-nest}
For a renormalizable quadratic map, a principal nest is a nested sequence
\[
 V_0\supset V_1\supset V_2\supset\cdots
\]
of critical puzzle domains adapted to the renormalization tower.  Its
principal-nest annuli are
\[
 A_n^{\mathrm{pn}}:=V_n\setminus\overline{V_{n+1}}.
\]
\end{definition}

\begin{definition}[A priori bounds]
\label{def:apriori-bounds}
A tower has uniform a priori bounds if there exists $\mu>0$ such that
\[
 \modulus(A_n^{\mathrm{pn}})\ge\mu
\qquad\text{for every }n.
\]
\end{definition}

\begin{lemma}[Uniform lower bounds imply divergence]
\label{lem:uniform-bound-divergence}
If $\modulus(A_n^{\mathrm{pn}})\ge\mu>0$ for every $n$, then
\[
 \sum_n\modulus(A_n^{\mathrm{pn}})=+\infty.
\]
\end{lemma}

\begin{proof}
The partial sums dominate $N\mu$ and therefore are unbounded.
\end{proof}

\begin{theorem}[Renormalization-to-puzzle shrinking]
\label{thm:renormalization-to-puzzle}
Suppose a satellite tower has a uniform principal-nest lower bound and that
the principal-nest annuli decompose into disjoint finite unions of the
ordinary puzzle annuli.  Then the ordinary puzzle-annulus moduli are not
summable, and the critical dynamical puzzle pieces shrink to $\{0\}$.
\end{theorem}

The first implication is a finite-sum comparison.  The second is the
Grötzsch--Yoccoz theorem from Chapter~1.  In Lean this chain is implemented
by
\begin{center}
\lean{PrincipalNestTarget.paraPuzzle_shrink_of_uniformConformalLowerBoundTarget}.
\end{center}

\section{Parameter-plane conclusion}

\begin{theorem}[Satellite a priori bounds imply pointwise MLC]
\label{thm:satellite-apriori-mlc}
Let $c\in\M$.  Assume:
\begin{enumerate}
\item a satellite renormalization tower at $c$;
\item a uniform principal-nest modulus lower bound;
\item connected parameter parapuzzle pieces at every depth.
\end{enumerate}
Then $\M$ is locally connected at $c$.
\end{theorem}

\begin{proof}
The modulus lower bound gives singleton intersection of the parameter
parapuzzles by Theorem~\ref{thm:renormalization-to-puzzle}.  Connectedness
and the neighbourhood-basis property then give local connectedness by
Theorem~\ref{thm:topological-transport}.  The Lean theorem
\lean{localConnectivity_of_satelliteAPrioriCertificate} packages this
implication.
\end{proof}

\section{Primitive, satellite, and neutral regimes}

The renormalization type is not cosmetic:
\begin{itemize}
\item primitive renormalization produces disjoint little Julia copies and
      is the natural setting for principal-nest modulus estimates;
\item satellite renormalization accumulates around roots of hyperbolic
      components and requires control of combinatorial rotation data;
\item neutral renormalization introduces Siegel disks, bubbles, hedgehogs,
      and neutral cascades, where ordinary expanding estimates degenerate.
\end{itemize}

Lyubich's bounded-type and a priori-bound machinery, Dudko--Lyubich's
satellite and Feigenbaum results, and the 2025--2026 Kahn/Dudko/Lim papers
all fit into this common contract.  Their analytic estimates differ, but
the formal target is the same: construct a tower certificate, prove
modulus control, realize connected parameter pieces, and invoke the
topological transport theorem.

\begin{openestimate}[Renormalization existence and estimates]
\label{est:renormalization}
For each literature class, define the exact combinatorial predicate,
construct the corresponding marked tower, formalize the relevant a priori
bounds, and prove that its parameter parapuzzles satisfy the hypotheses of
Theorem~\ref{thm:satellite-apriori-mlc}.
\end{openestimate}
